\section{Expanded out-of-family neuromodulation benchmark}
\label{sec:rethink-results}

The theory document motivates a substantially harder target than the earlier matched
complete-state benchmark: nonlinear high-dimensional dynamics, hidden modulators,
non-Gaussian and correlated innovations, calcium filtering and saturation, shared
artifact, partial observation, controlled operations, receptor kinetics, and
semi-Markov state paths. We therefore retained the earlier suite as a unit and
estimand-validation tier and executed a separate E1--E5 stress suite. No fitted method
received the hidden modulator or regime state.

The main run used 24 latent neurons, four hidden modulators, 12 training episodes, four
validation episodes, six test episodes, horizons 1, 4, and 16, and three independent
seeds. Seven fast methods ran in every cell. Gaussian and denoising neural networks,
an MDN, and feature-bilinear SBTG ran on a predeclared subset of 42 representative
stress cells. This produced 840 successful method--cell fits with no dropped failures.
A separate 32-neuron confirmation produced 672 successful fits for the full fast panel.

\subsection{Gaussian NLL is the empirical winner}

Table~\ref{tab:rethink-nll} reports average test negative log likelihood (lower is
better). Gaussian NLL ranks first in every scientific lane. This includes the
out-of-family nonlinear reservoir (E1), the paired measurement stress (E2), joint
noise/input/state/intervention shift (E3), stochastic release and receptor kinetics
(E4), and semi-Markov state dynamics (E5). The result persists at 32 neurons. Thus the
earlier Gaussian result was not only a consequence of giving a matched mechanistic
method the exact complete state, although the new generators remain bounded and the
available sample size still favors regularized low-complexity predictors.

\begin{table}[ht]
\centering
\caption{Mean test NLL in the 24-neuron expanded benchmark. Lower is better. Neural methods use their predeclared representative-cell panel and are not averaged over the same case mix as the fast methods.}
\label{tab:rethink-nll}
\begin{tabular}{@{}lrrrrr@{}}
\toprule
Method & E1 & E2 & E3 & E4 & E5 \\
\midrule
Gaussian NLL & 0.150 & -0.732 & -0.680 & -0.819 & -0.466 \\
Ridge VAR & 0.169 & -0.714 & -0.671 & -0.804 & -0.424 \\
Student-$t$ ridge & 0.255 & -0.645 & -0.600 & -0.728 & -0.331 \\
Gaussian MLP NLL & 0.203 & -0.229 & -0.628 & -0.663 & -0.386 \\
MDN & 0.206 & -0.217 & -0.608 & -0.634 & -0.393 \\
\bottomrule
\end{tabular}
\end{table}

Across all 96 fast-panel cells at 24 neurons, Gaussian NLL improves on ridge VAR by
0.0187 NLL units on average and wins 76.0\% of paired cells. At 32 neurons its advantage
increases to approximately 0.055 NLL units and it wins approximately 95\% of comparable
cells. With only three independent simulator seeds these are repeated-condition
summaries, not a precise population confidence interval. Seed-level Gaussian-NLL skill
over ridge is 0.001, 0.039, and 0.014 at 24 neurons, and 0.064, 0.031, and 0.070 at 32
neurons.

The nonlinear density models do not exploit a decisive advantage at this data scale.
Gaussian MLP, DSM-trained Gaussian MLP, and MDN remain valid but generally trail the
linear Gaussian likelihood. SBTG is a score-only model and is therefore excluded from
the NLL table rather than assigned a fictitious likelihood. Its held-out DSM risk is
finite across all 42 cells, but that native loss cannot establish predictive-law or
typed-channel superiority over normalized-density baselines.

\subsection{Unconstrained quadratic SID fails expected-risk and stability gates}

The quadratic SID median fit can appear reasonable, but its expected predictive risk is
not. At 24 neurons, the Hyv\"arinen and DSM variants pass their own numerical-validity
diagnostic in 90.6\% and 92.7\% of cells, respectively. At 32 neurons these rates fall to
82.3\% and 90.6\%. Rare crossings of the natural-variance boundary create extremely
large means or variances, so mean NLL and mean normalized RMSE become catastrophic even
when median NLL looks ordinary. This is not a cosmetic outlier issue: a probabilistic
model used for forecasting or likelihood comparison must control expected loss.

A post-benchmark diagnostic fit 432 additional SID models across 18 representative
E1--E5 cells, crossing six ridge values and four corruption scales. Ridge 1 with DSM
corruption $\sigma=1$ was the best validity-first setting: 100\% validity, median NLL
$-0.526$, mean NLL $-0.351$, and mean normalized RMSE 0.769. The matched ridge-VAR mean
NLL was $-0.435$. Ordinary tuning therefore can stabilize the implementation, but it
does not make SID win. Moreover, some settings marked formally valid still had extreme
RMSE, showing that the present diagnostic is weaker than the desired deployment gate.

The theoretical implication is constructive: retain the score-identified estimands,
but replace the unconstrained quadratic natural-variance map with a parameterization
that enforces positive, bounded-away-from-zero scale and controls tail behavior by
construction. A constrained log-scale model, calibrated fallback, or normalized
density backend may serve SID readouts without inheriting this failure mode.

\subsection{Paired intervention response favors direct regularized dynamics}

E3 and E4 include common-random-number control and operated trajectories. We evaluate
the difference between predicted operated and control outcomes against the true paired
difference. Averaged across active operations, ridge, heteroskedastic ridge, and
Student-$t$ ridge tie at response NRMSE 0.189; Gaussian NLL scores 0.194;
DSM-SID 0.202; Hyv\"arinen SID 0.210; and the zero-response baseline 0.293. The methods
capture part of the response through the evolving observed history, but all attenuate
its magnitude. No current SID method has an intervention-transfer advantage.

These response estimates are reduced-form. E4 uses stochastic release, saturating
occupancy, adaptation, and receptor attenuation, but good fluorescence forecasting does
not identify a unique ligand--receptor mechanism. Likewise, E5 establishes predictive
performance under modulator-dependent semi-Markov switching; a dedicated HMM/HSMM
state-recovery tournament remains necessary before claiming state or dwell-law recovery.

\subsection{Updated decision}

The expanded benchmark changes the priority ordering.

\begin{enumerate}
  \item Use Gaussian NLL and ridge as mandatory baselines and current practical defaults.
  \item Stop treating unconstrained quadratic SID as a deployable predictive law.
  \item Preserve finite typed SID contrasts as the main theoretical contribution, but
        compute them from a stable normalized density, direct ratio/classifier, or a
        positivity-constrained score backend.
  \item Continue diffusion only in a declared home field---multimodal, implicit, or
        sample-only conditional laws---and require it to beat MDN/flow/direct-ratio
        baselines there.
  \item Add explicit HMM/HSMM and direct response-kernel baselines before elevating the
        E5 path-law or E4 mechanistic claims.
\end{enumerate}
